% Chantier 95 -- Masked Mandates -- version IEEE Open Journal of the Computer Society (OJ-CS)
% Gabarit officiel IEEE Open Journals (ieeetj.cls, template-selector.ieee.org, telecharge le 12/11/2026).
% Build: pdflatex main-ojcs && bibtex main-ojcs && pdflatex main-ojcs && pdflatex main-ojcs (Overleaf : pdfLaTeX)
% numbers.tex, tab_*.tex and fig_*.pdf are generated by code/gen_tex_v21.py from the results files -- do not edit.
\documentclass{ieeetj}
\usepackage{cite}
\usepackage{iftex}
\ifXeTeX % local preview only (Overleaf compiles with pdfLaTeX and Times)
  \usepackage{fontspec}\setmainfont{texgyretermes}[Extension=.otf,UprightFont=*-regular,BoldFont=*-bold,ItalicFont=*-italic,BoldItalicFont=*-bolditalic]
\fi
\AtBeginDocument{\definecolor{tmlcncolor}{cmyk}{0.93,0.59,0.15,0.02}\definecolor{NavyBlue}{RGB}{0,86,125}}
\def\OJlogo{\vspace{-4pt}$<$Society logo(s) and publication title will appear here.$>$}
\def\seclogo{\vspace{10pt}$<$Society logo(s) and publication title will appear here.$>$}
\usepackage{amsmath,amssymb,amsthm}
\usepackage{booktabs}
\usepackage{array}
\usepackage{graphicx}
\usepackage{url}
\usepackage{xcolor}
\usepackage{multirow}
\usepackage{tikz}
\usepackage{hyperref}
\hypersetup{hidelinks=true}
\newtheorem{definition}{Definition}
\newtheorem{theorem}{Theorem}
\newtheorem{lemma}{Lemma}
\newcommand{\sys}{\textsc{Masked Mandates}}
\newcommand{\ok}{\checkmark}
\newcommand{\no}{$\times$}
\newcommand{\na}{--}
\newcommand{\pa}{$\circ$}
\newcommand{\Allow}{\mathsf{Allow}}
\newcommand{\Live}{\mathsf{Live}}
\newcommand{\Adv}{\mathsf{Adv}}
\newcommand{\negl}{\mathsf{negl}}
\newcommand{\ctx}{\mathsf{ctx}}
\newcommand{\nym}{\mathit{nym}}
\newtheorem{assumption}{Assumption}
\def\authorrefmark#1{\ensuremath{^{\textbf{#1}}}}
\makeatletter\renewcommand{\p@subsection}{\thesection-}\makeatother % refs to subsections read e.g. VII-D
% generated by gen_tex_v21.py from results/ -- do not edit
\newcommand{\AblBbsPres}{4}
\newcommand{\AblDyAdd}{165}
\newcommand{\AblHideCost}{0.5}
\newcommand{\AblNmAdd}{58}
\newcommand{\AblRangeAdd}{163}
\newcommand{\AblRevAdd}{165}
\newcommand{\AblRevPres}{168}
\newcommand{\BenchN}{100}
\newcommand{\CampSpread}{0.5\,\%}
\newcommand{\CoalFiveName}{51\,\%}
\newcommand{\CpuMac}{Apple M5 Pro}
\newcommand{\CpuZen}{AMD EPYC 9534}
\newcommand{\DecJoinCandDay}{5.5}
\newcommand{\DecJoinCandDayTenK}{44.8}
\newcommand{\DecJoinCandHour}{1.2}
\newcommand{\DecJoinCandHourTenK}{3.7}
\newcommand{\DecJoinPrecDay}{44\,\%}
\newcommand{\DecJoinPrecDayTenK}{23\,\%}
\newcommand{\DecJoinPrecHour}{83\,\%}
\newcommand{\DecJoinPrecHourTenK}{50\,\%}
\newcommand{\DelIssuerMs}{423}
\newcommand{\DelMsgHiS}{5.23}
\newcommand{\DelMsgLoMs}{1{,}293}
\newcommand{\DelParentS}{2.28}
\newcommand{\DelegLatMedS}{2.44}
\newcommand{\DepthPresMax}{390}
\newcommand{\DepthPresMin}{390}
\newcommand{\EBaseMaxHttp}{6.4}
\newcommand{\ECalls}{200}
\newcommand{\EChalHttp}{3.6}
\newcommand{\EFourCalls}{83{,}764}
\newcommand{\EFourCanonCalls}{83{,}764}
\newcommand{\EFourCfgs}{6}
\newcommand{\EFourCpu}{Intel~Xeon Platinum 8462Y+}
\newcommand{\EFourDelegS}{2.46}
\newcommand{\EFourDelegShareMax}{22\,\%}
\newcommand{\EFourDeltaDelayS}{1.2}
\newcommand{\EFourDeltaDelaySparseS}{1.2}
\newcommand{\EFourDeltaRej}{0}
\newcommand{\EFourDwLlmHonest}{2{,}147}
\newcommand{\EFourDwLlmSent}{629}
\newcommand{\EFourDwQuotaSent}{48}
\newcommand{\EFourDwRevSent}{48}
\newcommand{\EFourDwSent}{96}
\newcommand{\EFourEqAll}{14}
\newcommand{\EFourEqPairs}{15}
\newcommand{\EFourEqTsev}{12}
\newcommand{\EFourHonestCalls}{37{,}795}
\newcommand{\EFourIpcMs}{0.14}
\newcommand{\EFourJwtShareMax}{0.5\,\%}
\newcommand{\EFourJwtSlowMax}{1.00}
\newcommand{\EFourLatMed}{120}
\newcommand{\EFourLlmFixSuccMax}{95}
\newcommand{\EFourLlmFixSuccMin}{89}
\newcommand{\EFourMMEq}{5}
\newcommand{\EFourMMEqTsev}{4}
\newcommand{\EFourMMMaxDiff}{0.6}
\newcommand{\EFourMMMaxDiffTsev}{0.9}
\newcommand{\EFourMarginPP}{5}
\newcommand{\EFourMinHolm}{1.00}
\newcommand{\EFourMinP}{0.12}
\newcommand{\EFourModels}{5}
\newcommand{\EFourMtPres}{121}
\newcommand{\EFourMtVer}{63}
\newcommand{\EFourNTests}{30}
\newcommand{\EFourNoneSafeMax}{0.7}
\newcommand{\EFourNotEqTsev}{Granite-3.3-8B MM+escrow, gpt-oss-20b MM, gpt-oss-20b MM+escrow}
\newcommand{\EFourNudgeMax}{37}
\newcommand{\EFourNudgeMin}{0}
\newcommand{\EFourOrWzeroRej}{17\,\%}
\newcommand{\EFourPresMs}{353}
\newcommand{\EFourProbeCases}{63}
\newcommand{\EFourQuotaAtt}{3{,}735}
\newcommand{\EFourQuotaMax}{99}
\newcommand{\EFourQuotaMin}{2.1}
\newcommand{\EFourQuotaNone}{1{,}454}
\newcommand{\EFourRepDetModel}{Qwen2.5-72B}
\newcommand{\EFourRepMax}{7}
\newcommand{\EFourRepPP}{10}
\newcommand{\EFourRepPairs}{72}
\newcommand{\EFourRetryCalls}{696}
\newcommand{\EFourRetryDelayMs}{245}
\newcommand{\EFourRetryN}{192}
\newcommand{\EFourRetryOk}{192}
\newcommand{\EFourRetryRej}{0}
\newcommand{\EFourRevAtt}{1{,}954}
\newcommand{\EFourRevNone}{926}
\newcommand{\EFourRuns}{14{,}400}
\newcommand{\EFourSafeMax}{100}
\newcommand{\EFourSafeMin}{17}
\newcommand{\EFourSdk}{1.30.0}
\newcommand{\EFourShBaseMax}{94}
\newcommand{\EFourShBaseMin}{88}
\newcommand{\EFourShDelayTwoS}{4.2}
\newcommand{\EFourShDelayZeroMs}{239}
\newcommand{\EFourShEvery}{2}
\newcommand{\EFourShRejMax}{19\,\%}
\newcommand{\EFourShRejMin}{17\,\%}
\newcommand{\EFourShSuccTwoMax}{96}
\newcommand{\EFourShSuccTwoMin}{88}
\newcommand{\EFourShSuccZeroMax}{65}
\newcommand{\EFourShSuccZeroMin}{52}
\newcommand{\EFourShTasks}{120}
\newcommand{\EFourShareMax}{59\,\%}
\newcommand{\EFourShareMin}{35\,\%}
\newcommand{\EFourSlowMax}{2.32}
\newcommand{\EFourSlowMin}{1.49}
\newcommand{\EFourSlowNoLatMax}{3.12}
\newcommand{\EFourTasks}{120}
\newcommand{\EFourUnitsAll}{360}
\newcommand{\EFourUnitsTsev}{216}
\newcommand{\EFourVerMs}{95}
\newcommand{\EFourVllm}{0.33.0}
\newcommand{\EFourWtwoDelayS}{4.2}
\newcommand{\EFourWtwoDelaySparseS}{120}
\newcommand{\EMMBytes}{18{,}115}
\newcommand{\EMMDelegS}{4.01}
\newcommand{\EMMPrep}{291}
\newcommand{\EMMTotalDelay}{422}
\newcommand{\EMMTotalHi}{379}
\newcommand{\EMMTotalHttp}{377}
\newcommand{\EMMTotalLo}{376}
\newcommand{\EMMVer}{78}
\newcommand{\EMcpSdk}{1.30.0}
\newcommand{\EPy}{3.13.16}
\newcommand{\ERuns}{3}
\newcommand{\EStatelessAdd}{26}
\newcommand{\EStatelessAddJwt}{25}
\newcommand{\EWarm}{15}
\newcommand{\EscExtraB}{500}
\newcommand{\EtoeNoneHttp}{3.0}
\newcommand{\GapMedS}{0.78}
\newcommand{\GapModels}{5}
\newcommand{\GapN}{8{,}527}
\newcommand{\GapNinetyNineS}{5.6}
\newcommand{\GapNinetyS}{2.1}
\newcommand{\GapRows}{1{,}776}
\newcommand{\GapSource}{E4 v2.1 \sys{} runs}
\newcommand{\GitCommit}{96a694d14ca4}
\newcommand{\IssExactA}{71\,\%}
\newcommand{\IssExactAPdOne}{71\,\%}
\newcommand{\IssExactAPdOneTenK}{19\,\%}
\newcommand{\IssExactATenK}{19\,\%}
\newcommand{\IssExactI}{75\,\%}
\newcommand{\IssExactIPdOne}{87\,\%}
\newcommand{\IssExactIPdOneTenK}{33\,\%}
\newcommand{\IssExactITenK}{22\,\%}
\newcommand{\IssLagMedS}{1.06}
\newcommand{\IssLagNinetyNineS}{1.75}
\newcommand{\IssNDeleg}{556}
\newcommand{\IssNameInc}{1.1}
\newcommand{\IssNameIncHi}{1.4}
\newcommand{\IssNameIncHiPdOne}{5.2}
\newcommand{\IssNameIncHiPdOneTenK}{8.2}
\newcommand{\IssNameIncHiTenK}{1.7}
\newcommand{\IssNameIncLo}{0.9}
\newcommand{\IssNameIncLoPdOne}{3.7}
\newcommand{\IssNameIncLoPdOneTenK}{3.9}
\newcommand{\IssNameIncLoTenK}{0.8}
\newcommand{\IssNameIncPdOne}{4.3}
\newcommand{\IssNameIncPdOneTenK}{5.7}
\newcommand{\IssNameIncTenK}{1.2}
\newcommand{\IssNameWith}{70\,\%}
\newcommand{\IssNameWithPdOne}{73\,\%}
\newcommand{\IssNameWithPdOneTenK}{47\,\%}
\newcommand{\IssNameWithTenK}{43\,\%}
\newcommand{\IssPDeleg}{31\,\%}
\newcommand{\IssSessInc}{1.1}
\newcommand{\IssSessIncLoadHi}{0.0}
\newcommand{\IssSessIncLoadHiRate}{104}
\newcommand{\IssSessIncPdOne}{3.7}
\newcommand{\IssSessIncPdOneTenK}{4.8}
\newcommand{\IssSessIncTenK}{1.0}
\newcommand{\JitDSixHundred}{0.1\,\%}
\newcommand{\JitDSixty}{0.1\,\%}
\newcommand{\JitDTen}{0.0\,\%}
\newcommand{\JitNameSixHundred}{34\,\%}
\newcommand{\JitNameSixty}{37\,\%}
\newcommand{\JitNameTen}{50\,\%}
\newcommand{\LinkCalls}{179{,}802}
\newcommand{\LinkDayReid}{13\,\%}
\newcommand{\LinkDays}{30}
\newcommand{\LinkMMMinReid}{0.5\,\%}
\newcommand{\LinkMMMinSensMax}{2.0\,\%}
\newcommand{\LinkMMReid}{6.7\,\%}
\newcommand{\LinkMMReidHi}{10\,\%}
\newcommand{\LinkMMReidLo}{3.3\,\%}
\newcommand{\LinkMMUpper}{43\,\%}
\newcommand{\LinkPertoolReid}{94\,\%}
\newcommand{\LinkPtejOracle}{83\,\%}
\newcommand{\LinkPtejReid}{24\,\%}
\newcommand{\LinkSeeds}{5}
\newcommand{\LinkTools}{50}
\newcommand{\LoadRateMax}{104}
\newcommand{\LoadRateMin}{0.35}
\newcommand{\LoadUsersMax}{100{,}000}
\newcommand{\MsgBMax}{14{,}688}
\newcommand{\MsgBMin}{13{,}152}
\newcommand{\MsgHi}{65}
\newcommand{\MsgLo}{17}
\newcommand{\MsgPresMax}{392}
\newcommand{\MsgPresMin}{390}
\newcommand{\NAttMDel}{7}
\newcommand{\NAttMP}{13}
\newcommand{\NAttMPReach}{12}
\newcommand{\NAttProver}{12}
\newcommand{\NAttReach}{84}
\newcommand{\NAttacks}{96}
\newcommand{\NBitFlips}{2{,}000}
\newcommand{\NChecks}{37}
\newcommand{\NMsgs}{29}
\newcommand{\NameDay}{70\,\%}
\newcommand{\NameDecDay}{20\,\%}
\newcommand{\NameDecDayAdapt}{34\,\%}
\newcommand{\NameDecDayAdaptHi}{46\,\%}
\newcommand{\NameDecDayAdaptLo}{25\,\%}
\newcommand{\NameDecDayAdaptSess}{21\,\%}
\newcommand{\NameDecDayAdaptSessTenK}{15\,\%}
\newcommand{\NameDecDayAdaptTenK}{18\,\%}
\newcommand{\NameDecDayAdaptTenKHi}{22\,\%}
\newcommand{\NameDecDayAdaptTenKLo}{14\,\%}
\newcommand{\NameDecDayEvery}{61\,\%}
\newcommand{\NameDecDayEveryHi}{72\,\%}
\newcommand{\NameDecDayEveryLo}{50\,\%}
\newcommand{\NameDecDayFree}{17\,\%}
\newcommand{\NameDecDayFreeHi}{23\,\%}
\newcommand{\NameDecDayFreeLo}{11\,\%}
\newcommand{\NameDecDayFreeSess}{13\,\%}
\newcommand{\NameDecDayFreeSessTenK}{12\,\%}
\newcommand{\NameDecDayFreeTenK}{13\,\%}
\newcommand{\NameDecDayFreeTenKHi}{16\,\%}
\newcommand{\NameDecDayFreeTenKLo}{9.2\,\%}
\newcommand{\NameDecDayFreeTools}{20\,\%}
\newcommand{\NameDecDayFreeToolsHi}{25\,\%}
\newcommand{\NameDecDayFreeToolsLo}{14\,\%}
\newcommand{\NameDecDayFreeToolsTenK}{15\,\%}
\newcommand{\NameDecDayFreeToolsTenKHi}{19\,\%}
\newcommand{\NameDecDayFreeToolsTenKLo}{11\,\%}
\newcommand{\NameDecDayHi}{25\,\%}
\newcommand{\NameDecDayLo}{14\,\%}
\newcommand{\NameDecDayMin}{15\,\%}
\newcommand{\NameDecDaySess}{14\,\%}
\newcommand{\NameDecDayTenK}{15\,\%}
\newcommand{\NameDecDayTenKHi}{19\,\%}
\newcommand{\NameDecDayTenKLo}{11\,\%}
\newcommand{\NameDecHour}{22\,\%}
\newcommand{\NameDecHourAdapt}{44\,\%}
\newcommand{\NameDecHourAdaptHi}{60\,\%}
\newcommand{\NameDecHourAdaptLo}{30\,\%}
\newcommand{\NameDecHourAdaptSess}{29\,\%}
\newcommand{\NameDecHourAdaptSessTenK}{20\,\%}
\newcommand{\NameDecHourAdaptTenK}{24\,\%}
\newcommand{\NameDecHourAdaptTenKHi}{32\,\%}
\newcommand{\NameDecHourAdaptTenKLo}{17\,\%}
\newcommand{\NameDecHourEvery}{61\,\%}
\newcommand{\NameDecHourEveryHi}{72\,\%}
\newcommand{\NameDecHourEveryLo}{51\,\%}
\newcommand{\NameDecHourFree}{20\,\%}
\newcommand{\NameDecHourFreeHi}{25\,\%}
\newcommand{\NameDecHourFreeLo}{14\,\%}
\newcommand{\NameDecHourFreeSess}{14\,\%}
\newcommand{\NameDecHourFreeSessTenK}{13\,\%}
\newcommand{\NameDecHourFreeTenK}{14\,\%}
\newcommand{\NameDecHourFreeTenKHi}{17\,\%}
\newcommand{\NameDecHourFreeTenKLo}{10\,\%}
\newcommand{\NameDecHourFreeTools}{20\,\%}
\newcommand{\NameDecHourFreeToolsHi}{26\,\%}
\newcommand{\NameDecHourFreeToolsLo}{14\,\%}
\newcommand{\NameDecHourFreeToolsTenK}{15\,\%}
\newcommand{\NameDecHourFreeToolsTenKHi}{19\,\%}
\newcommand{\NameDecHourFreeToolsTenKLo}{11\,\%}
\newcommand{\NameDecHourHi}{26\,\%}
\newcommand{\NameDecHourLo}{15\,\%}
\newcommand{\NameDecHourMin}{15\,\%}
\newcommand{\NameDecHourSess}{14\,\%}
\newcommand{\NameDecHourTenK}{16\,\%}
\newcommand{\NameDecHourTenKHi}{20\,\%}
\newcommand{\NameDecHourTenKLo}{11\,\%}
\newcommand{\NameLoadHi}{30\,\%}
\newcommand{\NameLoadHiRate}{31}
\newcommand{\NameLoadHiUsers}{30{,}000}
\newcommand{\NameLoadLo}{71\,\%}
\newcommand{\NameLoadLoRate}{0.35}
\newcommand{\NameMM}{68\,\%}
\newcommand{\NameMMEvery}{89\,\%}
\newcommand{\NameMMHi}{79\,\%}
\newcommand{\NameMMLo}{54\,\%}
\newcommand{\NameMMMin}{56\,\%}
\newcommand{\NameMMSess}{54\,\%}
\newcommand{\NameMMSessPrec}{92\,\%}
\newcommand{\NameMMTenK}{41\,\%}
\newcommand{\NameMMTenKHi}{52\,\%}
\newcommand{\NameMMTenKLo}{31\,\%}
\newcommand{\NameSessLoadHi}{25\,\%}
\newcommand{\NameSessLoadHiRate}{104}
\newcommand{\NameSessLoadHiUsers}{100{,}000}
\newcommand{\NameSessLoadLo}{58\,\%}
\newcommand{\OmegaKB}{4.8}
\newcommand{\PresB}{13{,}536}
\newcommand{\PresKB}{13.5}
\newcommand{\PresMacMs}{194}
\newcommand{\PresMs}{390}
\newcommand{\ProfDecDay}{2.3\,\%}
\newcommand{\ProfDecDayHi}{4.3\,\%}
\newcommand{\ProfDecDayLo}{1.3\,\%}
\newcommand{\ProfDecHour}{1.9\,\%}
\newcommand{\ProfDecHourHi}{3.6\,\%}
\newcommand{\ProfDecHourLo}{1.1\,\%}
\newcommand{\ProfLoadHi}{0\,\%}
\newcommand{\ProfLoadHiRate}{31}
\newcommand{\ProfLoadLo}{6.2\,\%}
\newcommand{\ProfLoadLoRate}{0.35}
\newcommand{\ProfLoadMax}{10\,\%}
\newcommand{\PubParamsKB}{0.5}
\newcommand{\RevMax}{$10^{6}$}
\newcommand{\RevPresMax}{391}
\newcommand{\RevPresMin}{390}
\newcommand{\RevRefreshMs}{5.3}
\newcommand{\RevWitMs}{35}
\newcommand{\ThrBadMean}{19}
\newcommand{\ThrBadMs}{13}
\newcommand{\ThrBadPninety}{57}
\newcommand{\ThrGarbageUs}{1}
\newcommand{\ThrMpCores}{64}
\newcommand{\ThrMpEff}{91\,\%}
\newcommand{\ThrMpRate}{556}
\newcommand{\ThrPerCore}{10}
\newcommand{\ThrVerMs}{104}
\newcommand{\VerBytesMs}{104}
\newcommand{\VerMacMs}{42}
\newcommand{\VerMs}{85}
\newcommand{\Warmup}{5}

\makeatletter\let\inputtab\@@input\makeatother % primitive input: safe inside tabular

\begin{document}
\receiveddate{XX Month, XXXX}
\reviseddate{XX Month, XXXX}
\accepteddate{XX Month, XXXX}
\publisheddate{XX Month, XXXX}
\currentdate{XX Month, XXXX}
\doiinfo{XXXX.2026.1234567}

\markboth{}{Ameur et al.}

\title{Masked Mandates: Unlinkable, Attenuable and Accountable Delegation for LLM Agents}

\author{Yulliwas Ameur\authorrefmark{1,2}, Samia Bouzefrane\authorrefmark{2}, and Soumya Banerjee\authorrefmark{2}}
\affil{Efrei Research Lab, Universit\'e Paris-Panth\'eon-Assas, Paris, France}
\affil{C\'edric, Conservatoire national des arts et m\'etiers (Cnam), Paris, France}
\corresp{Corresponding author: Yulliwas Ameur (email: yulliwas.ameur@efrei.fr; ORCID: 0000-0003-1435-2982).}
\authornote{ORCID: Y.~Ameur 0000-0003-1435-2982; S.~Bouzefrane 0000-0002-0979-1289; S.~Banerjee 0000-0001-8766-8473.}

\begin{abstract}
Agents built on large language models (LLMs) call tools and delegate sub-tasks, yet current delegation tokens show every tool whom the agent serves, so colluding tools can profile a person. With \sys{}, an agent proves, for one tool call, only that it holds an unrevoked mandate from a recognized issuer allowing this action within a per-call cap, an expiry and a rate limit shared by its whole delegation tree. Each call carries one zero-knowledge proof bound to it and a pseudonym with a hidden counter; revoking a mandate revokes its descendants; a two-of-three escrow can identify the principal after abuse. We prove six properties, including unlinkability against a corrupted issuer, assuming a consistent log and simulation-extractable proofs. A presentation takes \PresMs{}~ms to produce and \VerMs{}~ms to verify on one core. With five open LLMs as agents, it slows tasks \EFourSlowMin{} to \EFourSlowMax{} times and enforces revocation and quotas; under a shared issuer, \EFourShRejMin{}--\EFourShRejMax{} of honest calls fail unless tools briefly accept the previous accumulator value. Timing still links: in simulation, once one tool learns identities, colluding tools name \NameMM{} of calls at low load and \NameLoadHi{} at high load.
\end{abstract}
\begin{IEEEkeywords}
Anonymous credentials, autonomous agents, delegation, Model Context Protocol, privacy, rate limiting, revocation, zero-knowledge proofs.
\end{IEEEkeywords}

\maketitle

% Sections 1--6 of the revised OJ-CS manuscript (input by main-ojcs.tex), split into three files.
% Sections 1--6 of the revised OJ-CS manuscript (input by main-ojcs.tex). Hand-written; numbers come from numbers.tex.
\section{Introduction}
\IEEEPARstart{A}{gents} act on behalf of people. To book a trip, an agent built on a large language model (LLM) may query an airline, a hotel and a payment service, and may hand the hotel search to a sub-agent. Each call needs authorization: the tool must know that the caller may perform \emph{this} action. Current designs answer with a token that carries identifiers and the delegation chain:
signed intent tokens with per-agent keys~\cite{goswami2025ajwt}; JSON Web Tokens (JWT) or Biscuit tokens for multi-hop delegation over the Model Context Protocol (MCP) and Agent2Agent (A2A)~\cite{prakash2026aip}; OAuth token exchange and transaction tokens~\cite{rfc8693,tulshibagwale2026txntokens}; persistent decentralized identities~\cite{zou2026intersage,garzon2026didvc}.
These tokens prove authority, but disclose the principal and give every tool a stable identifier to log. One mandate may drive dozens of tools in a day; each tool sees a slice, and pooled together the slices form a profile.

\paragraph*{Problem} Can an agent prove that it may make one specific tool call, without revealing who mandated it and without being linkable across tools by the tokens it shows? The tool should still be able to rate-limit the whole delegation tree, the principal to revoke a subtree, and a committee to lift anonymity after abuse.

\paragraph*{What exists} Each property exists somewhere:
delegatable anonymous credentials give anonymous, attenuable delegation, some with chained revocation~\cite{belenkiy2009dac,camenisch2017cdd,griffy2024dac,mir2024tdac,xie2025fdaac}; updatable credentials re-issue hidden attributes through an online issuer~\cite{blomer2019uacs,hartung2017bbaplus}; $k$-times anonymous authentication, Clone Wars, anonymous counting tokens, Privacy Pass and anonymous rate-limited credentials (ARC) rate-limit anonymously~\cite{teranishi2004ktaa,au2006bbsplus,camenisch2006clonewars,benhamouda2023act,rfc9576,yun2026arccrypto}; group signatures give traceability with non-frameability~\cite{bellare2005dynamic}.
None of these targets a delegation tree of agents whose every call must be bound to one tool invocation and counted against the root's quota. The agent-specific schemes either carry identities~\cite{goswami2025ajwt,prakash2026aip,zou2026intersage} or do not address delegation~\cite{liu2026bcm,llambi2026cva,zhou2026capbinding}.

\paragraph*{Contributions} Our contribution is a careful composition for agents, not a new primitive.
\begin{enumerate}
\item A deployment and threat model for agent delegation through a shared domain issuer, and six game-based properties (Sec.~\ref{sec:model}): unforgeability with attenuation and chained revocation, a rate limit shared by a whole delegation tree, invocation binding, unlinkability with an explicit leakage function, traceability, and non-frameability against a corrupted issuer that publishes a consistent log.
\item \sys{} (Sec.~\ref{sec:construction}), which combines BBS+ mandates~\cite{au2006bbsplus,tessaro2023bbs} delegated by blind re-issuance; one Dodis--Yampolskiy pseudonym per call, with a hidden counter taken from a signed interval, so that sub-agents get disjoint shares of the root's quota~\cite{dodis2005vrf,camenisch2006clonewars}; a per-root tag bound by the principal's signature; a VB accumulator with hash-derived parameters for chained revocation~\cite{vitto2022vb}; and a 2-of-3 escrow that opens to the principal and the mandate. We prove the properties under stated assumptions in the supplementary material.
\item An open Rust implementation, Python bindings and an MCP binding with tool-issued nonces (Sec.~\ref{sec:impl}).
\item An evaluation (Sec.~\ref{sec:eval}): an attack taxonomy with malicious provers and fuzzing; costs on two platforms, an ablation and multi-core throughput; an MCP comparison with five token baselines; a collusion study with identifier, content, timing and identified-tool attackers; and agent runs with five open LLMs, including delegation, revocation, binding quotas and a shared issuer.
\end{enumerate}

\section{Background and Related Work}\label{sec:related}
\subsection{Agent identity and delegation}
Agentic JWT ties each agent action to the user's intent and carries a chain of downstream agents with per-agent proof-of-possession keys~\cite{goswami2025ajwt}. AIP uses a signed JWT for single-hop delegation and a Biscuit token~\cite{biscuit} for multi-hop delegation across MCP and A2A, verified locally~\cite{prakash2026aip}. OAuth token exchange, rich authorization requests and transaction tokens propagate identity and scope along call chains~\cite{rfc8693,rfc9396,tulshibagwale2026txntokens}; the MCP authorization specification builds on OAuth~\cite{mcpauth}. Other agent-specific proposals:
InterSAGE, a positioning paper, proposes attenuated agent credentials with subtree revocation~\cite{zou2026intersage}; Rodriguez Garzon et al.\ give agents decentralized identifiers (DIDs) and verifiable credentials~\cite{garzon2026didvc}; KITA requires threshold, action-bound signatures~\cite{zheng2026kita}; South et al.\ argue for authenticated delegation to agents~\cite{south2025icml}.
Google's Agent Payments Protocol (AP2) carries signed ``mandates'' as verifiable credentials over A2A and MCP and shows them to merchants as evidence~\cite{ap2-2025}. The OpenID Foundation~\cite{south2025idmgmt} and the IETF WIMSE working group~\cite{kasselman2026aims} survey identity and delegation for agents; the WIMSE draft recommends minimizing personal data in tokens. These systems target authority and accountability; none aims at unlinkability. B-CM authenticates agents unlinkably~\cite{liu2026bcm}, and Zhou binds capabilities with BBS+ certificates~\cite{zhou2026capbinding}; neither addresses delegation between agents or per-call binding. A zero-knowledge authorization proof for agents uses a public agent identifier~\cite{llambi2026cva}.

\subsection{Anonymous credentials and delegation}
\paragraph*{BBS-type signatures and related credentials} BBS-type signatures~\cite{bbs2004,au2006bbsplus,camenisch2016bbsplus} allow efficient proofs with selective disclosure. Their security was revisited recently~\cite{tessaro2023bbs,chairattana2025bbs}, issuance can be distributed among several signers~\cite{doerner2023threshold}, per-verifier pseudonyms are specified in an IRTF CFRG draft~\cite{kalos2026pseudonyms}, and a W3C cryptosuite brings BBS to verifiable credentials~\cite{w3c2026bbs}. Related credential systems:
Coconut thresholdizes issuance~\cite{sonnino2019coconut}; issuer-hiding credentials hide which issuer signed~\cite{bobolz2021ihabc,connolly2022spseq}; zk-creds builds credentials from general-purpose zero-knowledge proofs~\cite{rosenberg2023zkcreds}; keyed-verification credentials serve Signal groups~\cite{chase2020signal}.

\paragraph*{Delegation} Delegatable credentials~\cite{belenkiy2009dac,camenisch2017cdd} became practical with mercurial signatures~\cite{griffy2024dac}, gained threshold and fine-grained delegation with chained revocation~\cite{mir2024tdac,xie2025fdaac}, and are built in recent preprints from JWTs with recursive SNARKs~\cite{rometsch2026dac} and in the keyed-verification setting~\cite{mirzamohammadi2026dkvac}. Updatable anonymous credentials let an online issuer update hidden attributes under a proven relation~\cite{blomer2019uacs}, as BBA+ does for balances~\cite{hartung2017bbaplus}. Our delegation is such an update: the parent proves a relation between its signed attributes and the child's committed ones, and the issuer blind-signs. Unlike delegatable credentials, it needs the online issuer, which in exchange enforces a per-root delegation limit and chooses identifiers.

\subsection{Anonymous rate limiting, tokens and blocklisting}
$k$-TAA~\cite{teranishi2004ktaa,au2006bbsplus} and periodic $n$-times anonymous authentication~\cite{camenisch2006clonewars} limit each credential to $k$ uses per period, with pseudonyms derived from a secret and a counter. We reuse this design; what is new is splitting the counter interval across a delegation tree. Anonymous counting tokens~\cite{benhamouda2023act}, Privacy Pass~\cite{rfc9576,rfc9577,rfc9578,hanff2025privacypass} and ARC~\cite{yun2026arccrypto} give unlinkable, rate-limited tokens bound to an origin challenge, but carry no attributes and no delegation. SD-JWT with batch issuance~\cite{rfc9901,openid4vci} gives one-time credentials that tools cannot link, but the issuer can. Anonymous blocklisting~\cite{tsang2011nymble,tsang2007blac,rosenberg2022snarkblock} blocks misbehaving users without identifying them, an alternative to escrow.

\subsection{Positioning}
Table~\ref{tab:position} summarizes the properties as stated in each work; the supplementary material and the artifact give the evidence for the partial and inferred cells. ``Not addressed'' does not mean ``impossible'': any Fiat--Shamir credential can bind a call by hashing it, and a delegatable credential could add $k$-TAA-style per-verifier limits. What is new here is one hidden counter space split across a delegation tree, chained revocation through ancestor identifiers in a BBS+ setting, and their integration with escrow. \sys{} has all the listed properties, but it is also one of the few schemes that need an online issuer to delegate (last column), and it costs far more per call than token schemes (Sec.~\ref{sec:eval}). Sec.~\ref{sec:discussion} says when each choice fits.

\begin{table}[t]\centering
\caption{Properties as described in each work: \ok{} provided, \pa{} partially, \no{} stated as not provided, \na{} not addressed, ? not verifiable from the open text, n/a not applicable. Unl.: unlinkable. Del.: delegation with attenuation. Bnd.: proof bound to one call or a verifier challenge (Privacy Pass, ARC, ACT: challenge, not the call). RL: anonymous per-verifier rate limit. Sub.: revoking a credential revokes its descendants. Tr.: escrow or tracing. Off.: delegation without an online issuer. $^*$: our inference, not a claim of the work.}\label{tab:position}
\scriptsize\setlength{\tabcolsep}{2.2pt}
\begin{tabular}{@{}lccccccc@{}}\toprule
System & Unl. & Del. & Bnd. & RL & Sub. & Tr. & Off.\\\midrule
Agentic JWT~\cite{goswami2025ajwt} & \no & \pa & \pa & \na & \na & \ok$^5$ & \no\\
AIP~\cite{prakash2026aip} & \no & \ok & \no$^4$ & \na & \no & \ok$^5$ & \ok\\
InterSAGE$^1$~\cite{zou2026intersage} & \no & \ok & \na & \na & \ok & \ok$^5$ & \no\\
DID/VC agents~\cite{garzon2026didvc} & \no & \na & \na & \na & \na & \ok$^5$ & \na\\
KITA~\cite{zheng2026kita} & \na & \na & \ok & \na & \na & n/a & \na\\
AP2~\cite{ap2-2025} & \no & \pa$^*$ & \ok & \na & \na & \ok$^5$ & \na\\
Verif.\ authorization~\cite{llambi2026cva} & \no & \no & \ok & \na & \na & \ok$^5$ & \na\\
B-CM$^2$~\cite{liu2026bcm} & \ok & \na & \na & \na & \na & ? & \na\\
CDD17~\cite{camenisch2017cdd} & \ok & \ok & \na & \na & \na & \na & \ok\\
DAC (mercurial)$^3$~\cite{griffy2024dac} & \ok & \pa & \na & \na & \ok & \pa & \ok\\
TDAC$^2$~\cite{mir2024tdac} & \ok & \ok & \na & \na & \na & \na & ?\\
FDAAC-CR~\cite{xie2025fdaac} & \ok & \ok & \na & \na & \ok & \na & \ok\\
UACS~\cite{blomer2019uacs} & \ok & \pa$^*$ & \na & \na & \na & \pa & \no\\
BBA+~\cite{hartung2017bbaplus} & \ok & \na & \na & \na & \na & \pa & n/a\\
DAC from JWTs~\cite{rometsch2026dac} & \ok & \pa$^*$ & \na & \na & \na & \na & \ok$^*$\\
Delegatable KVAC~\cite{mirzamohammadi2026dkvac} & \ok & \ok & \na & \na & \na & \na & ?\\
BBS per-verifier pseud.~\cite{kalos2026pseudonyms} & \pa & \na & \pa$^*$ & \na & \na & \na & n/a\\
Clone Wars~\cite{camenisch2006clonewars} & \ok & \na & \na & \ok & \na & \pa & n/a\\
Dynamic $k$-TAA~\cite{au2006bbsplus} & \ok & \na & \na & \ok & \na & \pa & n/a\\
BLAC/SNARKBlock~\cite{tsang2007blac,rosenberg2022snarkblock} & \ok & \na & \na & \na & \na & \pa & n/a\\
ACT (counting tokens)~\cite{benhamouda2023act} & \ok & \na & \pa & \ok & \na & \na & n/a\\
Privacy Pass~\cite{rfc9577,hanff2025privacypass} & \ok & \na & \pa & \pa & \na & \na & n/a\\
ARC~\cite{yun2026arccrypto} & \ok & \na & \pa & \ok & \na & \na & n/a\\
SD-JWT, batch~\cite{rfc9901} & \pa & \na & \pa & \na & \na & n/a & n/a\\
Coconut~\cite{sonnino2019coconut} & \ok & \na & \na & \pa & \na & \na & n/a\\
\sys{} & \ok & \ok & \ok & \ok & \ok & \ok & \no\\\bottomrule
\multicolumn{8}{@{}p{0.97\columnwidth}@{}}{\scriptsize $^1$Positioning paper, no construction. $^2$Closed access; title and abstract only. $^3$Assumes an honest revocation authority that can link chains. $^4$Its invocation-bound tokens do not prevent replay within their lifetime. $^5$Trivially traceable: the identity is in clear.}
\end{tabular}
\end{table}


% Sections 3 (deployment, threat model, properties) of the revised OJ-CS manuscript.
\section{Deployment, Threat Model and Properties}\label{sec:model}
\subsection{Deployment}
The parties are as follows (a figure is in the supplementary material).
\begin{itemize}
\item \emph{Issuer} $\mathcal{I}$: a domain authorization service shared by many principals, for example the authorization server of an enterprise or of a federation of agent platforms. It must not be the agent host, which sees the conversations anyway. It publishes its public key, every accumulator value and every root-tag binding in an append-only log; wallets accept nothing that is not in the log, so the issuer cannot tag wallets with private values.
\item \emph{Principal} $P$: registers once a long-term key $pk_P=g^{u_P}$, kept in its \emph{wallet}, by a self-signed entry in the log, which it monitors; the registry is therefore not in the issuer's hands. To obtain a root mandate, $P$ authenticates to $\mathcal{I}$ as usual and signs a fresh per-root tag $T=g^{t}$ with $u_P$.
\item \emph{Wallet}: keeps the tree's secrets and produces presentations when an \emph{agent} asks for one. The agent is an LLM-driven process that never sees the secrets. The wallet enforces a policy: the time must be close to its clock, the amount must equal the amount in the call and be in the stated currency, and an optional allow-list of tools applies.
\item \emph{Sub-agents}: each receives a delegated mandate with its own disjoint interval of counters. Cryptographic attenuation matters when a sub-agent runs its own wallet (threat A2), which then holds the tree's secrets and can link its calls; with a shared wallet, attenuation could be wallet policy (E4's delegations ran in one process).
\item \emph{Tools} (MCP servers, A2A agents): verify presentations with public data, their clock, their nonces and a seen-set of pseudonyms. Each tool publishes a fixed descriptor: its action class, the JSON path of the amount, the currency, and whether it requires escrow.
\item \emph{Escrow committee}: three members, which we recommend be independent of the issuer and of the platform, hold shares of a decryption key from a distributed key generation (DKG).
\end{itemize}

\paragraph*{Defaults and anonymity set} Epochs last one hour; roots are long-lived (say one per principal and month). The issuer limits roots per principal and bounds each root's interval by a public $Q$ ($2^{20}$ here): root issuance is the Sybil boundary. The anonymity set of a presentation is the set of principals of the same issuer whose mandates allow the same call in the same epoch, against the same accumulator value. A tool learns that the caller is authorized by $\mathcal{I}$, the action class, and the values it set itself (time, amount). It does not learn the principal, the mandate or its position in the tree; it can learn the remaining quota only by probing until refused.

\paragraph*{Where it fits} The scheme pays off when one issuer serves many principals that tools must not tell apart: a platform or federation that imposes data minimization on third-party tools, or tools that sell per-call access; for an enterprise issuer, tools still learn the employer. Tools that need the user's name (payment, delivery) gain nothing for their own calls, and should be called in a separate session at another time, carrying no value echoed from the task (such as a quoted price), which in E3 cuts naming from \NameMM{} to \NameDecHourFree{} (\NameDecHourEvery{} if every task needs such a tool); in our travel example, \sys{} protects the searches and the sub-agent's activity, not the payment. The supplementary material details these settings and the GDPR roles.

\paragraph*{Running example} Alice's wallet obtains a root mandate for classes \{search, pay\}, a cap of 1{,}000~EUR per call and counters $[0,100)$ per tool and hour.
Her travel agent delegates hotel booking to a sub-agent with class \{search\}, a cap of 300~EUR and counters $[80,100)$; the hotel tool sees only a 13.5-kB proof that some mandate of its issuer allows this exact call, and a fresh pseudonym; Alice's tree makes at most 100 calls per hour at each tool, however many sub-agents it has; revoking the sub-agent or the root stops it or the whole tree everywhere; after abuse, two escrow members can name Alice and the mandate.

\subsection{Threat model}
A1: \emph{Linking}: any coalition of tools, possibly with the issuer and one escrow member, tries to link presentations to each other or to a principal. A2: \emph{Misbehaving agent}: a compromised or prompt-injected agent~\cite{greshake2023injection}, or a malicious wallet, tries to act beyond its mandate, give a sub-agent more than it holds, or forge proofs. A3: \emph{Replay} to another tool, call, epoch or nonce. A4: \emph{Rate abuse}, including minting sub-agents or inventing epochs to multiply quotas. A5: \emph{Use after revocation} of a mandate or an ancestor. A6: \emph{Framing}: the issuer, tools and one escrow member try to attribute a presentation to an innocent principal.

\paragraph*{Outside the cryptographic model} Unlinkability concerns the authorization material only. Call content (measured in E3), network addresses (use a relay such as Oblivious HTTP~\cite{rfc9458}), MCP session state (use one stateless session per call; E2 measures its cost), per-client OAuth tokens~\cite{mcpauth}, and the issuer's view of identified root issuance and delegation times can still identify callers; the supplementary material lists the countermeasures. The agent host and LLM provider see the conversation and are trusted by the principal.
Prompt injection itself is out of scope; \sys{} bounds what an injected agent can do. We assume TLS and authenticated tool identities.

\subsection{Properties}\label{sec:defs}
Table~\ref{tab:props} states the six properties informally, with the coalition each one resists and its assumptions; the supplementary material gives the games and proofs. All proofs use random oracles. Extraction uses rewinding with sequential issuer sessions, so a deployment must serialize issuance. The composite Fiat--Shamir proof is assumed simulation-extractable, for its $\Sigma$-protocol part~\cite{faust2012fsnm} and for the range arguments, which need knowledge soundness rather than plain soundness~\cite{attema2022fsmulti,ganesh2022bpse}. Unlinkability with an opening oracle additionally needs a chosen-ciphertext-secure escrow, which holds for signed ElGamal in the algebraic group model~\cite{fuchsbauer2020schnorr}.

\begin{table*}[t]\centering\caption{Security properties (games and proofs in the supplementary material). q-SDH: strong Diffie--Hellman; DL: discrete logarithm; q-DDHI: decisional Diffie--Hellman inversion; ZK: zero knowledge; SE: simulation extractability.}\label{tab:props}
\scriptsize
\begin{tabular}{@{}p{0.1\textwidth}p{0.42\textwidth}p{0.2\textwidth}p{0.2\textwidth}@{}}\toprule
Property & Statement & Adversary & Assumptions\\\midrule
UNF & Every accepted call extracts to a mandate that descends by attenuation from a root the issuer signed, allows the call, and is live (no revoked ancestor). & all agents, tools, escrow; honest issuer & q-SDH, DL, SE (incl.\ range arguments)\\
RL & Per root, tool and epoch of the tool's clock: at most $k_0-\ell_0$ accepted calls by the whole tree, and as many delegations per issuer epoch. & as UNF & as UNF\\
INV & A presentation cannot be used for another call, tool, epoch or nonce. & issuer, tools, other principals & q-DL, ZK, SE\\
UNL & Presentations by two honest mandates are indistinguishable when the leakage $\Lambda$ (allowed, live, accumulator value, valid witness, quota left) is equal; many challenges, each advancing both mandates; with or without opening oracle. & issuer, all tools, one escrow member & ZK, DDH, $(Q{+}1)$-DDHI, consistent log (CCA: signed ElGamal)\\
TRACE & Two escrow members open every accepted call to its principal and mandate. & agents, tools, one escrow member; honest issuer & as UNF, Schnorr\\
NF & No coalition can make a presentation open to an honest principal. & issuer, tools, one escrow member & DL, Schnorr, ZK, SE, consistent log\\\bottomrule
\end{tabular}
\end{table*}


% Sections 4--6 (construction, security, implementation) of the revised OJ-CS manuscript.
\section{Construction}\label{sec:construction}
\subsection{Building blocks}
We work in BLS12-381.
BBS+ signs a vector of messages; a holder proves knowledge of a signature while revealing any subset, and Pedersen commitments allow blind issuance~\cite{au2006bbsplus,camenisch2016bbsplus,pedersen1991}; a VB accumulator holds revoked identifiers; a holder proves that a hidden value is not a member~\cite{vitto2022vb}; Bulletproofs++ prove that a committed value lies in a range~\cite{eagen2024bpp}.
All statements of one operation form a single Fiat--Shamir proof, in which witness equalities force the same hidden value across statements~\cite{dockcrypto}. All bases and proving keys are derived by hashing public labels. This matters: an issuer that chose the accumulator's proving key could strip the randomization of non-membership proofs and recover hidden identifiers.

\subsection{Mandate layout}
A mandate is a BBS+ signature on $L=13+n$ messages ($n=16$ classes, $L=\NMsgs$): $s$, $\mathit{id}$, $\mathit{anc}_{1..4}$, $\delta$, the scope bits $b_{1..n}$, $\beta$, $\ell$, $k$, $\tau$, $t$ and a key version. All mandates of a tree share $s$ and $t$. The issuer chooses $\mathit{id}$; $s$ and $t$ are never shown to it.

\subsection{Algorithms}
\paragraph*{IssueRoot} The wallet draws $s,t$ and sends $C=h_0^{r}h_s^{s}h_t^{t}$, $T=g^{t}$, a Schnorr signature by $u_P$ on (label, $T$), and a proof that $C$ commits to $\log_g T$, bound to an issuer nonce. $\mathcal{I}$ authenticates $P$ and checks the proof, the signature, that $T$ is new, the layout of the signed vector, the bounds $\beta,k,\tau<2^{40}$ and $k-\ell\le Q$, and its per-principal root limit. It then draws a fresh non-zero $\mathit{id}$, blind-signs, and logs the binding. The wallet fixes every clear attribute ($\ell_0,k_0,\beta,\tau$, scope) in its request and checks that the returned signature is on exactly those values.

\paragraph*{Present} An invocation $\iota=(\theta,\mathsf{call},c,\mathit{amt},\mathit{cur},e,\mathit{tm},n_\theta)$ consists of the tool's public identity $\theta$, the canonical call, its class, amount and currency, the tool's epoch and time, and a nonce issued by the tool. For $\iota$, the wallet picks the next unused counter $j$ in its own part of $[\ell,k)$ for $(\theta,e)$ and computes $\nym=B_{\theta,e}^{1/(s+j+1)}$ with $B_{\theta,e}=H(\theta,e)$; the base depends only on public tool data and the epoch, never on the tree's secrets $s$ and $t$. It proves, with Fiat--Shamir context $\ctx=\mathrm{SHA\mbox{-}256}(\iota)$ and against the accumulator value $V$ it holds:
\begin{align*}
\mathsf{PoK}\{&\text{signature on } m \text{ with } b_c=1,\\
&\nym^{s+j+1}=B_{\theta,e},\ \ell\le j<k,\ \beta\ge\mathit{amt},\ \tau\ge \mathit{tm},\\
&\mathit{id},\mathit{anc}_{1..4}\notin\mathsf{Acc}_V,\\
&[\,E=(g^{r_1},g^{t}X^{r_1},g^{r_2},g^{\mathit{id}}X^{r_2})\,]\}.
\end{align*}
$\ell\le j<k$ is proved by two ranges, $j-\ell+1\in[1,2^{40})$ and $k-j-1\in[0,2^{40})$. The escrow part $E$ is included when the tool requires it. The output is $(V,\nym,[E],\pi)$; $j$ stays hidden.

\paragraph*{Verify} The tool performs these checks in order:
(1) rebuilds $\iota$ from the raw request, its identity, its clock and the nonce it issued; the nonce is consumed; (2) checks that $e$ is its current epoch and the time is within tolerance; (3) checks that $V$ is the current accumulator value or one of the $W$ previous ones; (4) reserves $(e,\nym)$ in its seen-set; (5) verifies $\pi$, and releases the reservation on failure.
All failures return the same error.

\paragraph*{Delegate} The issuer sends a nonce. The parent chooses the child's attributes, normally with a sub-interval $[\ell',k')$ disjoint from its own part and from its other children. It commits to all child messages except $\mathit{id}'$ and the version ($C_p$), computes a delegation pseudonym $\nym_D=D_e^{1/(s+j+1)}$ for a counter in its own part, and proves, with context $H(\text{nonce}\|C_p)$:
its signature and the non-membership of its five identifiers; by witness equalities: $s'=s$, $t'=t$, $\mathit{anc}'=(\mathit{id},\mathit{anc}_1,\mathit{anc}_2,\mathit{anc}_3)$; $\delta'=\delta+1\le4$; $\beta'\le\beta$, $\tau'\le\tau$, $\ell'\ge\ell$, $k'\le k$, each as a difference in $[1,2^{40})$, together with child values in $[0,2^{40})$; for each class, $b_c-b'_c+1\in[1,2]$; the pseudonym.
The issuer checks that $e$ is its own current epoch, verifies, rejects a seen $(e,\nym_D)$, draws a fresh $\mathit{id}'$ and blind-signs. The parent unblinds and hands the mandate to the child. The issuer learns $\mathit{id}'$, the epoch and the time of the delegation, and nothing else.

The ranges on the child values make attenuation hold over the integers; with differences alone, a child value could wrap around $q$. The scope rule needs no bit proof: a presentation reveals $b_c=1$ exactly, and by induction no descendant of a $0$ (or negative) value can reach $1$.

\paragraph*{Revoke} The log accepts a revocation of $X\neq0$ only with an anonymous proof made with the credential of the live mandate $X$ itself, which only its holder and the parent that delegated it keep (so no descendant revokes an ancestor), or with signatures of two escrow members; identifier $0$, the padding for missing ancestors, is never revocable. The issuer then appends $X$ to the accumulator and publishes the new value with its update data and a spot check (witnesses, under both values, for a point hashed from the update), which anyone can verify; holders refresh their five witnesses from public data, and a wallet that cannot refresh stops presenting. Because each mandate carries its ancestors' identifiers, revoking one identifier stops a subtree. Presentations in flight survive unrelated revocations through a time-bounded window: a tool accepts the previous value only for $\Delta$ seconds after its successor reaches it, so revocation takes effect after $\Delta$ plus propagation, and with batched updates at most two values coexist. A count-based window or a wallet retry are alternatives (E4); the retry links the two attempts at one tool.

\paragraph*{Escrow and Judge} Committee members hold shares of $x_{\mathcal{E}}$ from the DKG of Gennaro et al.~\cite{gennaro2007dkg}. Member $i$ returns $E_1^{x_i}$ and $E_3^{x_i}$ with Chaum--Pedersen proofs~\cite{chaum1992dleq}. Two shares give $g^{t}$ and $g^{\mathit{id}}$. $\mathsf{Judge}$ maps $g^{t}$ to a logged binding, re-verifies the principal's signature on it, and maps $g^{\mathit{id}}$ to the issued mandate.

\section{Security}\label{sec:security}
The supplementary material proves the six properties of Table~\ref{tab:props}. Four points deserve mention.
\begin{itemize}
\item \emph{Arithmetic}: the bound check for $[\ell,h)$ proves $m-\ell$ and $h-1-m$ in $[0,2^{64})$; with $\ell,h<2^{41}<2^{65}<q$, every attenuation, interval and counter statement holds over the integers.
\item \emph{Quota}: within one tool identity $\theta$ (origin and tool name) and epoch, the pseudonym is an injective function of the counter, and the counter lies in the root's interval. Epochs are the tools' and the issuer's, so they cannot be invented. Sub-intervals give children shares of the quota that add up to the root's without coordination. A parent that hands out overlapping intervals only causes rejections inside its own tree.
\item \emph{Issuer}: unlinkability holds against a corrupted issuer that publishes a consistent log, and non-frameability as long as registrations are in that log rather than in the issuer's database. Because the log accepts revocations only from a mandate's holder or parent or from two committee members, the issuer cannot revoke a root to watch which traffic stops; it can still hand one wallet an invalid initial witness, which that wallet detects (tampered update data fail the public spot check).
\item \emph{Trees}: one tree shares $s$ and $t$. A delegate that runs its own wallet can link that tree's calls, and its actions open to the principal. It never receives $u_P$, so it cannot obtain roots or affect the principal's other trees.
\item \emph{Concrete security}: the unlinkability reduction loses $n_h^2Q$ for $n_h$ honest trees and rests on $(Q+1)$-DDHI (about $2^{117.5}$ under Cheon's attack for $Q=2^{20}$), so it gives $2^{80}$ only for a few hundred honest trees per issuer key: deployments should keep $Q$ small and rotate the key; this limits the proof, it is not an attack.
\end{itemize}

\section{Implementation}\label{sec:impl}
{\sloppy The prototype is a Rust crate on \texttt{proof\_system}~0.34, \texttt{bbs\_plus}~0.25, \texttt{vb\_accumulator}~0.29 and \texttt{bulletproofs\_plus\_plus}~0.9~\cite{dockcrypto}, over BLS12-381. A PyO3 module exposes two roles:
an \texttt{Authority}, which plays the issuer, the registry, the escrow and the wallets of the test harness, and allocates disjoint intervals; a \texttt{Verifier}, built from public data only, which issues nonces and checks the epoch and the time.
The MCP binding carries $\{V,\nym,[E],\pi\}$ in the request's \texttt{\_meta} field~\cite{mcpspec}. Each server exposes a \texttt{challenge} tool that returns a single-use nonce and the server's time. The server rebuilds $\iota$ from its origin and tool name, the RFC~8785 canonical JSON of the raw call~\cite{rfc8785}, and its descriptor (class, amount, currency). The \texttt{challenge} tool is called by the client, not offered to the language model. The binding, which we have not yet proposed upstream, is outside the core MCP authorization specification, which requires an OAuth bearer token on every HTTP request and forbids servers to accept other tokens~\cite{mcpauth}; it fits the specification's optional authorization extensions, alongside a per-server access token issued to the domain and shared by all its principals, which identifies the domain, not the user (supplementary material). The same envelope fits A2A message metadata~\cite{a2aspec}, which we did not evaluate. The implementation was written with an AI assistant from the construction above and reviewed by the first author (see Acknowledgment); it has not been audited independently.\par}


% Evaluation section of the revised OJ-CS manuscript (input by main-ojcs.tex). Every number is a macro from numbers.tex.
\section{Evaluation}\label{sec:eval}
\paragraph*{Setup} Cryptographic costs (E1) are measured with one thread (\texttt{RAYON\_NUM\_THREADS=1}) on two platforms. The server platform is a core of an \CpuZen{} (MesoNET Zen): each of three independent campaigns ran with its own seed on its own exclusive node, with the process pinned to one core, and all three used one binary built from commit \texttt{\GitCommit}, whose hash is recorded with the results. A uniformly slower first run of one campaign was excluded by a rule fixed before its rerun and rerun on another node (supplementary material; both runs are in the artifact). For comparison, the laptop platform is one core of an \CpuMac{}; the laptop was not idle during the measurements, and its load is recorded with the results. Each point has \Warmup{} warm-up runs and at least 30 measured runs (\BenchN{} for presentation and verification). We report medians with 95\,\% confidence intervals from a hierarchical bootstrap, which resamples campaigns and then runs. Campaign medians differ by at most \CampSpread{} for every operation. A script generates every number, table and figure from the result files.

\subsection{Attack suite}
All \NAttacks{} attacks are rejected and all \NChecks{} correctness checks pass; \NAttReach{} attacks reach a verifier (a tool, or anyone checking published data), the issuer, the revocation log or Judge, and \NAttProver{} are refused by the honest prover library, which shows nothing about verifier soundness. The supplementary material classifies them by threat and rejecting party. The checks cover honest presentations at depths 0 to 4 with and without escrow, every member of a tree using its own part of the interval, a strict tool, the grace window, 2-of-3 openings, and tools re-deriving every proof parameter from public labels.

\paragraph*{Malicious provers and delegators} \NAttMPReach{} of the \NAttMP{} malicious-prover attacks bypass every client-side check and are rejected by the tool or the issuer, which rebuild the statement from their own inputs. Tools reject proofs for a weaker statement (amount or time bound 0), for another revealed class, or with a quota statement detached from the signed interval. For \NAttMDel{} malicious delegations, one per attenuation relation (child budget, interval bounds, expiry, scope bit or depth beyond the parent's, delegation counter beyond its quota), the prover's statement for that relation is re-bound to the parent's own value; the issuer, which rebuilds the statement, rejects each one, and a debug issuer that omits only the binding of that relation to the child's commitment accepts it (positive control). These tests show that every relation is bound to the child's committed values; the relation itself rests on the knowledge soundness of the range arguments, which a test cannot exercise.

\paragraph*{Deployment and framing attacks} The issuer rejects a delegation for another epoch, a root interval wider than its bound $Q$, a root beyond its per-principal limit, and a root request whose attribute layout is shifted by an extra scope entry (which would otherwise move every later field and evade $Q$ and tracing). Tools reject a presentation for another epoch, at a time far from their clock, with a nonce they never issued or already consumed, or against an accumulator value older than the grace window. The log rejects unauthorized, forged, replayed or stale revocations, a descendant revoking its ancestor, and identifier 0. Judge rejects a tag that the issuer binds to a principal without the principal's signature, and, with the registry taken from the log, a key that the issuer substitutes for the principal's. Finally, \NBitFlips{} random single-bit flips over the whole presentation are all rejected.

\subsection{Costs (E1)}
Table~\ref{tab:costs} gives each operation for the default layout (\NMsgs{} messages).
\begin{itemize}
\item \emph{Presentation}: \PresMs{}~ms to produce and \VerMs{}~ms to verify on the server core (\VerBytesMs{}~ms including deserialization and subgroup checks), and \PresMacMs{} and \VerMacMs{}~ms on the laptop core. A presentation is \PresB{} bytes; escrow adds \EscExtraB{} bytes.
\item \emph{Delegation}: \DelParentS{}~s for the parent and \DelIssuerMs{}~ms for the issuer, once per sub-agent.
\item \emph{Public data}: a tool needs \PubParamsKB{}~kB of public data. Everything else is derived from labels.
\end{itemize}

\begin{table}[t]\centering\caption{Costs, default layout, one core: median [95\,\% CI] on the server core, median on the laptop core (ms); bytes.}\label{tab:costs}
\footnotesize\setlength{\tabcolsep}{2pt}
\begin{tabular}{@{}lrrr@{}}\toprule Operation & Server & Laptop & Bytes\\\midrule
\inputtab tab_costs.tex \end{tabular}\end{table}

\paragraph*{Where the time goes} Adding the statement groups one by one (table in the supplementary material): The BBS+ proof alone costs \AblBbsPres{}~ms. A pseudonym with a counter checked against the signed interval adds \AblRevAdd{}~ms if the counter is revealed, as in $k$-TAA, and \AblDyAdd{}~ms if it is hidden: hiding the counter costs \AblHideCost{}~ms, spent in the two range proofs of the interval. The cap and expiry ranges add \AblRangeAdd{}~ms, and the five non-membership proofs for chained revocation \AblNmAdd{}~ms. The four range proofs dominate; each uses two 64-bit Bulletproofs++ arguments in the library we use, so aggregating them is the obvious optimization.

\paragraph*{Scaling} Presentation cost does not depend on the size of the mandate, its depth or the revocation list: \MsgPresMin{}--\MsgPresMax{}~ms from \MsgLo{} to \MsgHi{} signed messages, \DepthPresMin{}--\DepthPresMax{}~ms for depths 0 to 4, and \RevPresMin{}--\RevPresMax{}~ms up to \RevMax{} revoked entries (figure and table in the supplementary material). Delegation grows linearly with the number of classes (\DelMsgLoMs{}~ms at \MsgLo{} messages, \DelMsgHiS{}~s at \MsgHi{}); after 100 revocations, each holder refreshes its five witnesses in \RevRefreshMs{}~ms.

\paragraph*{Throughput} Deserializing and verifying a presentation takes \ThrVerMs{}~ms, so one core verifies about \ThrPerCore{} presentations per second. Independent single-threaded verifiers scale almost linearly: \ThrMpRate{} per second on \ThrMpCores{} cores, a parallel efficiency of \ThrMpEff{}. A presentation with one flipped bit is rejected after \ThrBadMs{}~ms, and random bytes of the right length are rejected by the parser in about \ThrGarbageUs{}~$\mu$s.

\subsection{End-to-end on MCP (E2)}
\paragraph*{Setup} One agent client calls three MCP servers (Python SDK \EMcpSdk{}, Python \EPy{}), with cryptography on one laptop thread. There are four transports: stdio, streamable HTTP on the loopback interface, HTTP through a relay that adds 10~ms in each direction, and stateless HTTP, with a new connection and MCP session for every call so that no session identifier links calls.
Each mode and transport makes \ECalls{} calls after \EWarm{} warm-up calls. The whole measurement was repeated \ERuns{} times with different seeds with the machine load recorded; we report the median of the run medians. Before each authorized call the agent fetches a single-use nonce and the tool's time from the tool (\EChalHttp{}~ms over HTTP). All modes use the same two-hop delegation (root, agent, sub-agent) and bind each call to the canonical JSON (with currency), the tool identity and the nonce, except where the mechanism cannot.

\paragraph*{Baselines} (1) A JWT chain with one delegation JWT per hop and a per-call proof-of-possession JWT, in the style of Agentic JWT and DPoP~\cite{goswami2025ajwt,rfc9449}. (2) A Biscuit token with two attenuation blocks and a per-call block~\cite{prakash2026aip}. (3) A macaroon with two delegation caveats and a per-call caveat~\cite{birgisson2014macaroons}; tools must share the root key. (4) Batch one-time tokens: the issuer checks the delegation chain in the clear and issues single-use signed tokens with a proof-of-possession key, as with SD-JWT batch issuance~\cite{rfc9901}. (5) RSA blind-signature tokens (RFC~9474) that commit to an ephemeral key which signs the call, as in Privacy Pass with a proof-of-possession binding~\cite{rfc9474,rfc9578}. We added a status list to the JWT chain and to batch tokens so that they can revoke. ARC and anonymous counting tokens are not implemented.

\begin{table*}[t]\centering\caption{End-to-end MCP call (E2) on a laptop that was not idle (load recorded), medians over 3 runs: agent preparation, tool verification and total per call in ms (streamable HTTP, +10~ms per direction, stateless HTTP); authorization bytes; delegation (two hops) or issuance cost per token; whether a modified call is rejected, revoking the agent stops the sub-agent, colluding tools can link calls; what the issuer learns.}\label{tab:e2e}
\footnotesize\setlength{\tabcolsep}{2.5pt}
\begin{tabular}{@{}lrrrrrrrcccl@{}}\toprule Mode & Prep. & Verify & Total & +10 ms & Stateless & Bytes & Deleg./issuance & Binds & Revokes & Tools link & Issuer learns\\\midrule
\inputtab tab_e2e.tex \multicolumn{12}{@{}l}{\scriptsize $^\dagger$ except by a 2-of-3 escrow opening. Blind signing in pure Python.}\end{tabular}\end{table*}

\paragraph*{Results} Table~\ref{tab:e2e} shows the price of unlinkability.
\emph{Cost}: a local HTTP call costs \EtoeNoneHttp{}~ms without authorization, and every baseline stays at or below \EBaseMaxHttp{}~ms. \sys{} takes \EMMPrep{}~ms on the agent and \EMMVer{}~ms on the tool, for \EMMTotalHttp{}~ms per call (\EMMTotalLo{}--\EMMTotalHi{} across runs) and \EMMTotalDelay{}~ms with the added delay. Its two-hop delegation costs \EMMDelegS{}~s once. On the wire a presentation exceeds its \PresB{} raw bytes (base64 in the JSON \texttt{\_meta} field). For scale, a $k$-TAA-style credential with a revealed counter costs \AblRevPres{}~ms to present (ablation); blind RSA issuance is pure Python. \emph{Stateless sessions}: a fresh MCP session per call, which removes the session identifier as a linking handle, adds about \EStatelessAdd{}~ms per call in every challenge-based mode. \emph{Controls}: every mode except ``none'' rejects replays to the same and to another tool, also with a fresh nonce. Without proof of possession, blind tokens would accept a modified call or currency. Macaroons and blind tokens still accept the sub-agent after its parent is revoked. \emph{What the issuer learns}: with batch tokens, the issuer sees every delegation chain and token identifier. With blind tokens, it sees the chain and how many tokens each root requests. With \sys{}, it sees only when delegations happen.
Only \sys{} also offers a tree-wide quota and, with escrow, accountability.

\subsection{Collusion study (E3)}\label{sec:link}
\paragraph*{What we simulate} E3 measures what channels other than the token still reveal. A simulator generates 1{,}000 and 10{,}000 users calling \LinkTools{} tools with Zipf popularity over \LinkDays{} days, which gives \LinkCalls{} calls at 1{,}000 users.
Each user has favorite tools, interests, a home city, a budget and a daily routine. A task is a burst of 2 to 6 calls; the gaps between calls are resampled from \GapN{} gaps of the \sys{} runs of E4 (median \GapMedS{}~s, 99th percentile \GapNinetyNineS{}~s), and timestamps have millisecond resolution. A share of tasks (\IssPDeleg{}, as in E4) starts by delegating to a sub-agent. Booking and payment calls carry personal fields (city, party size, preferences), search calls carry topics, and a payment repeats to the cent a price returned earlier in the task.
All tools pool their logs.

\paragraph*{Attackers} (a) Identifier matching. (b) Timing: links calls closer than a window, and also runs a session tracker. (c) Sessions, then 30-day clustering on call content. (d) As (c), plus tool sets and exact amounts. The \emph{oracle} assigns each session to the user with the most likely true profile (strong, not an upper bound). In the \emph{named} setting, one tool per task (payment or booking) legitimately learns the user's name, and the coalition attaches the rest of the session, and then the clusters, to that name.

\paragraph*{Metrics} \emph{Profile recovery} is the share of users whose largest cluster holds at least half of their calls with at least 80\,\% purity. \emph{Naming} is the share of calls attributed to the right named person. We run \LinkSeeds{} seeds per cell. All numbers come from a synthetic population: they show which channels link and how strongly, not the rates of a real deployment. ``Minimized'' runs strip personal fields from arguments.

\begin{table}[t]\centering\caption{Collusion of 50 tools (E3), \%, mean of 5 seeds (ranges in the supplementary material): profile recovery by attacker (d) and by the profile oracle at 1{,}000 users, (d) at 10{,}000 users, and naming of calls when one tool per task learns the user's identity, at 1{,}000 and 10{,}000 users.}\label{tab:link}
\footnotesize\setlength{\tabcolsep}{2.5pt}
\begin{tabular}{@{}lrrrrr@{}}\toprule & \multicolumn{3}{c}{Profile recovery} & \multicolumn{2}{c}{Naming}\\ Token format & (d) & Oracle & (d), $10^4$ & $10^3$ & $10^4$\\\midrule
\inputtab tab_link.tex \end{tabular}\end{table}

Table~\ref{tab:link} gives four findings.
\begin{enumerate}
\item \emph{Identifiers matter.} Per-tool pseudonyms behave like a principal identifier once tools pool their logs: \LinkPertoolReid{} of profiles recovered. A revealed counter is also a strong signal: \LinkPtejReid{} by (d) and \LinkPtejOracle{} by the oracle, against \LinkMMReid{} and \LinkMMUpper{} for \sys{}. From the tools' view, \sys{} and one-time tokens give identical logs, and they get identical numbers.
\item \emph{Content still links.} With full arguments, (d) recovers \LinkMMReid{} of profiles (\LinkMMReidLo{}--\LinkMMReidHi{} over seeds). With minimized arguments it recovers \LinkMMMinReid{}, and at most \LinkMMMinSensMax{} across the sensitivity settings listed in the artifact.
\item \emph{A tool that learns the identity names the session.} Timing links the calls of a task, so the coalition names \NameMM{} of \sys{} calls at 1{,}000 users and low load (\NameMMLo{}--\NameMMHi{} over seeds). The share rises to \NameMMEvery{} if every task touches such a tool. Minimizing arguments only brings it to \NameMMMin{}, and five colluding tools already name \CoalFiveName{}. Naming weakens as tasks overlap: from \NameLoadLo{} at \LoadRateMin{} tasks per minute across the coalition to \NameLoadHi{} at \NameLoadHiRate{} (\NameLoadHiUsers{} users) (figure in the supplementary material). The attacker attributes every call, so the rest are named wrong. An issuer in the coalition, which knows when delegations happen (roots are long-lived), adds only \IssNameInc{} points (\IssNameIncPdOne{} if every task delegates).
\item \emph{Delays.} Random per-call delays of up to 10~s defeat profile clustering (\JitDTen{}) but not naming (\JitNameSixHundred{} with up to 600~s).
\end{enumerate}
Calling the identity-learning tool in its own session, at a random time within the next hour, cuts naming to \NameDecHour{}; but an attacker that joins the payment to the earlier quote of the same amount brings it back to \NameDecHourAdapt{}, and only a call that carries no value echoed from the task keeps it at \NameDecHourFree{} (\NameDecHourFreeTenK{} at 10{,}000 users). It stays at \NameDecHourEvery{} (original attacker) when every task needs such a tool. Unlinkable authorization is thus necessary but not sufficient: calls must also carry little, and identity-learning calls must be decoupled from the task and carry no value it echoes.

\subsection{LLM agents in the loop (E4)}\label{sec:e4}
\paragraph*{Setup} A language model served by vLLM~\EFourVllm{} on one H100 GPU solves \EFourTasks{} tasks by calling four MCP servers over streamable HTTP (SDK \EFourSdk{}). There are five templates: (a) pay the cheapest flight and hotel under caps; (b) the same, with the hotel delegated to a sub-agent; (c) as (b) plus a transfer, with the sub-agent revoked after its first payment (correct outcome: a reported failure); (d) a search-heavy task; (e) a five-city survey delegated with a three-search quota (correct outcome: flight paid, partial failure reported).
Each task runs in four modes: none, a JWT chain with status list and per-tree counters at the tools, \sys{}, and \sys{} with escrow. It runs in \EFourCfgs{} configurations: temperature 0 with and without tool latency (lognormal, median \EFourLatMed{}~ms), a replication of the first with identical settings, and temperature 0.7 with three seeds. This gives \EFourRuns{} runs over \EFourModels{} open models and \EFourCalls{} tool calls. Tools are strict verifiers; they canonicalize the raw JSON-RPC arguments themselves, with the same JCS implementation and MCP SDK as the client. Authorization travels in \texttt{\_meta}, so prompts, tool outputs and seeds are identical across modes. When a model stops without finishing, the harness sends up to two ``please continue'' messages (\EFourNudgeMin{}--\EFourNudgeMax{}\,\% of runs, depending on the model, in every mode). Agent and tools run on single cores of an \EFourCpu{} (presentation \EFourPresMs{}~ms, verification \EFourVerMs{}~ms).

\begin{table}[t]\centering\caption{LLM agents (E4), six configurations pooled, \%: success on honest templates (a, b, d); correct outcome of revocation (c) and quota (e) templates; median slowdown of \sys{} (MM) over none (tool latency on); share of task time in MM authorization (temperature 0).}\label{tab:e4}
\scriptsize\setlength{\tabcolsep}{1.8pt}
\begin{tabular}{@{}lrrrrrrrrr@{}}\toprule & \multicolumn{3}{c}{Honest} & \multicolumn{2}{c}{Revoked (c)} & \multicolumn{2}{c}{Quota (e)} & & \\ Model & none & JWT & MM & JWT & MM & JWT & MM & Slow. & Share\\\midrule
\inputtab tab_e4.tex \end{tabular}\end{table}

\paragraph*{Results} Table~\ref{tab:e4} answers four questions.
\begin{itemize}
\item \emph{Does authorization change what honest agents achieve?} By construction the model sees the same transcript in every mode, so outcomes can differ only through rejections or serving noise. The substantive result is that, with a per-task issuer, no honest call was wrongly rejected (\EFourHonestCalls{} calls with \sys{}), and client and tool canonical strings never differed, over \EFourCalls{} calls and \EFourProbeCases{} adversarial argument shapes. The rest is serving noise: rerunning a configuration changes up to \EFourRepMax{} of \EFourRepPairs{} outcomes (only \EFourRepDetModel{} is deterministic), so we use a pragmatic margin of $\pm$\EFourMarginPP{} points, half that churn, fixed after observing it. Equivalence means that the 90\,\% cluster-bootstrap interval (unit: task) lies within the margin, i.e.\ two one-sided tests at 5\,\%; per-configuration results are in the supplementary material. With a permutation test that respects the clustering of units in tasks, no difference to ``none'' on the honest templates is significant (smallest $p=\EFourMinP{}$; Holm over \EFourNTests{} tests). \sys{} is equivalent to ``none'' for \EFourMMEqTsev{} of 5 models at temperature 0.7 (not gpt-oss-20b) and for all 5 when the four distinct configurations are pooled (largest difference \EFourMMMaxDiff{} points); over all modes, \EFourEqTsev{} and \EFourEqAll{} of \EFourEqPairs{} mode--model pairs are equivalent.
\item \emph{Does it stop what it should?} All \EFourRevAtt{} calls by revoked sub-agents and all \EFourQuotaAtt{} calls beyond a sub-agent's quota were rejected under JWT and \sys{}. Without authorization, all \EFourRevNone{} and \EFourQuotaNone{} were accepted. Under \sys{} the honest wallet refused to present. A dishonest wallet that kept presenting after revocation (old accumulator value, or stale witness) and beyond its quota (reused counter) had all \EFourDwSent{} (scripted) and \EFourDwLlmSent{} (five models) such presentations rejected by the tools, with the uniform error. Whether the agent then reports the failure depends on the model: \EFourSafeMin{} to \EFourSafeMax{}\,\% correct outcomes for (c), \EFourQuotaMin{} to \EFourQuotaMax{}\,\% for (e). The small models often claim success after a rejection, with JWT as with \sys{}.
\item \emph{What does it cost?} With tool latency, \sys{} slows honest tasks by a median factor of \EFourSlowMin{} to \EFourSlowMax{} (up to \EFourSlowNoLatMax{} without tool latency), and JWT by at most \EFourJwtSlowMax{}. Authorization takes \EFourShareMin{} to \EFourShareMax{} of task time (JWT: at most \EFourJwtShareMax{}). Per-call proofs dominate; the sub-agent's delegation takes \EFourDelegS{}~s. The share falls with slower models or tools; a wallet in a separate process adds only \EFourIpcMs{}~ms per call.
\item \emph{What does a shared issuer change?} One issuer serves all tasks of a worker and revokes an unrelated mandate every \EFourShEvery{}~s. A presentation built just before a revocation and verified just after it is then rejected: with LLM agents, \EFourShRejMin{} to \EFourShRejMax{} of honest calls, and success falls to \EFourShSuccZeroMin{}--\EFourShSuccZeroMax{}\,\% (per-task issuer: \EFourShBaseMin{}--\EFourShBaseMax{}\,\%). A window of two accumulator values removes the rejections but leaves a revoked mandate usable for two updates (\EFourWtwoDelayS{}~s here, \EFourWtwoDelaySparseS{}~s at one revocation per minute). A wallet retry with the same counter and a fresh nonce rejects none of \EFourRetryCalls{} honest scripted calls, with revocation effective after \EFourRetryDelayMs{}~ms; a time-bounded window accepting the previous value for 1~s rejects none and bounds the delay to \EFourDeltaDelayS{}~s whatever the rate. Both delays exclude propagation from the log, and both remedies assume batched revocations (supplementary material). With Qwen2.5-72B and gpt-oss-20b, both reject no honest call and restore success to \EFourLlmFixSuccMin{}--\EFourLlmFixSuccMax{}\,\%.
\end{itemize}


% Discussion and conclusion of the revised OJ-CS manuscript (input by main-ojcs.tex).
\section{Discussion}\label{sec:discussion}
\paragraph*{When to use what} If the principal trusts the issuer with its activity and only tools must not link calls, batch one-time tokens are far cheaper (E2). If the issuer must not learn which calls a principal makes either, blind tokens with a proof-of-possession key suffice for flat authorization. \sys{} is for deployments that also need, against a curious issuer, delegation with attenuation, a quota shared by a tree, subtree revocation and accountability. Without the tree-wide quota, a revealed counter as in $k$-TAA would do; E1 shows that hiding it costs almost nothing once the interval check is there.

\paragraph*{Online issuer} Delegation through the issuer costs a round trip, makes the issuer a point of availability and tells it when delegations happen; in exchange it enforces the delegation limit and chooses identifiers. The consistent log limits it: it can withhold, visibly, but not tag. Delegatable credentials delegate offline (Table~\ref{tab:position}) but do not rate-limit a tree.

\paragraph*{What the rate limit covers} The quota is per tool name and epoch, so it multiplies with a server's tool names and doubles around an epoch boundary. $\beta$ caps each call, so spending is at most $(k_0-\ell_0)\beta_0$ per tool name and epoch (100 calls of 1{,}000~EUR per hour in our example), and $N$ times that across $N$ tools. Delegations per issuer epoch are bounded only by $k_0-\ell_0$, and a compromised parent can use its children's counters. A cumulative budget across tools, which AP2 mandates state for identified merchants~\cite{ap2-2025}, would need an online spend service, as with anonymous credit tokens~\cite{schlesinger2026act} (an expired draft, spent at the issuer), or a double-spend registry shared by tools; either learns when spends happen.

\paragraph*{Escrow governance} Opening is a decision about personal data. We recommend three committee members independent of the platform and of each other, a log commitment to every request and its justification before any share is released (opened once any investigation allows it), and escrow keys erased after the retention time. The supplementary material discusses the lawful basis of an opening (Art.~6), restricted notification (Art.~23) and, after \emph{EDPS v SRB}~\cite{cjeu2025srb}, when presentations are personal data for a tool.

\paragraph*{Denial of service and capacity} Verification scales with cores (\ThrMpRate{} per second on \ThrMpCores{} cores), and malformed inputs cost about \ThrGarbageUs{}~$\mu$s, but a well-formed invalid presentation costs a full verification. Tools should first require a cheap Privacy Pass token per nonce request~\cite{rfc9577}. Aggregating range proofs and precomputing first messages during generation remain unexplored.

\paragraph*{Limits} The construction is not post-quantum. Our proofs assume simulation-sound extraction for the composite proof, including simulation extractability of Bulletproofs++ inside it, which is assumed, and an algebraic-group-model argument for unlinkability with an opening oracle. The consistent log is assumed, not implemented; one in-process harness plays issuer, escrow and wallets, so delegation times exclude the round trip to the issuer, whose capacity (about \DelIssuerMs{}~ms per delegation) we did not study; the implementation has not been audited independently. The evaluation uses simulated tools and traffic, and open models served locally.

\section{Conclusion}
\sys{} lets an LLM agent prove exactly the authority it needs for one tool call: the mandate is valid, unrevoked and attenuated, the call is within the cap, and it is within the quota of the whole delegation tree. Assuming a consistent log and simulation-extractable proofs, it does not tell the tool whom the agent works for, and two of three escrow members can find the principal after abuse. Each call costs \PresMs{}~ms on the agent and \VerMs{}~ms on the tool on one server core, independently of delegation depth and of the revocation list. In agent runs with simulated tool latency this slows tasks by a factor of \EFourSlowMin{} to \EFourSlowMax{}; with an issuer shared by many tasks, \EFourShRejMin{}--\EFourShRejMax{} of honest calls fail unless tools accept the previous accumulator value for about a second. Unlinkable tokens are necessary but not sufficient: in our simulation, colluding tools recover \LinkMMReid{} of user profiles from call content (\LinkMMMinReid{} with minimized arguments, and \LinkMMUpper{} for an attacker that knows every user's profile), and when one tool learns the user's identity they name \NameMM{} of the calls through timing among 1{,}000 users at low load (\NameMMTenK{} and \NameLoadHi{} at 10{,}000 and 30{,}000 users).


\section*{Data availability}
Code, harnesses, simulators, result files with environment records and hashes, and the script that generates every number, table and figure are archived at Zenodo, concept DOI \href{https://doi.org/10.5281/zenodo.24148362}{10.5281/zenodo.24148362} (release 2.2.1, DOI \href{https://doi.org/10.5281/zenodo.24222783}{10.5281/zenodo.24222783}, and code tag v2.2.1 for this paper; 1.0.x hold the first re-implementation), under the MIT license (code) and CC BY 4.0 (results), with the evidence for Table~\ref{tab:position}.


\section*{Acknowledgment}
Y.~Ameur designed the scheme and wrote the manuscript, including the proofs, and the first prototype. S.~Banerjee contributed to the security definitions and proofs. S.~Bouzefrane supervised the work, discussed the scheme and reviewed the manuscript. After the first prototype was lost, the implementation released with this paper was re-built with an AI assistant, Claude (Anthropic), from the construction in Sec.~\ref{sec:construction}. Claude also wrote the attack suite, the benchmark scripts and the agent harnesses, carried out the revision of the implementation described in Sec.~\ref{sec:impl}, ran the measurements on the first author's laptop and on the MesoNET clusters, helped compare the related systems with the cited papers, and proposed edits to the text, each of which the first author approved. Y.~Ameur has reviewed the code, the proofs and the whole text; all authors approved the final manuscript and take full responsibility for its content. This work used the computing resources of MesoNET (project M26297, Hedy and Zen clusters). This work was supported by a French government grant managed by the Agence Nationale de la Recherche under the ``Investissements d'avenir'' program (reference ``ANR-21-ESRE-0051'').


\bibliographystyle{IEEEtran}
\bibliography{refs,refs2}
\end{document}
